\section{A word before the first Sunday}

What would you do if you knew the Lord was to be your guest before nightfall? In the Gospel a centurion comes to Jesus with a sick servant at home. He knows who needs help, he knows where to go, and he trusts that one word from Jesus can reach as far as his house. The Church reads his story on the first Monday of Advent (Matthew 8:5--11). At Mass, as the priest shows us the consecrated Host, we make the centurion's confession of unworthiness our own. Advent can begin there, with someone we love and something only Christ can do.

If you are preparing for Sunday Mass, turn first to the \hyperref[sec:pc-advent-sundays]{four Sundays of Advent}. Choose the letter printed beside the Sunday in your parish missal. Read the Gospel aloud once, slowly enough to hear the question or the promise. When you come to Mass, listen for the words on either side of the verse you remember. The person proclaiming the Gospel will often give back to you a part of the story you read too quickly at home.

If you have opened the book on a weekday, go to the \hyperref[sec:pc-advent-weekdays]{weekdays of Advent}. Through December 16, find the week and day. From December 17, find the date. Sunday has its own readings even when it falls on December 24. A solemnity or a feast may give your parish a different Gospel. The directions beside the readings will help you find your place. You can begin today. The Church will give you another First Sunday in due course.

A few minutes will do. Name before God the person who needs your prayer. Read the passage for the day, keep a little silence and pray the verse given there. If the house is full of children, let one child say what the people in the Gospel actually did. An answer as small as \emph{he went} or \emph{she listened} gives everyone something to take into the day. For a longer time of prayer, read the day's lesson as well as the Gospel. Take up the spiritual reading, then pray Morning or Evening Prayer from the Church's book. Someone sitting beside a hospital bed can read one sentence to a person too tired to speak. Silence after that sentence can be the prayer of them both.

\subsection{What we ask when we ask him to come}

The last petition of the Apocalypse is addressed to someone who has spoken first. Christ promises to come. Only then does the seer answer him (Revelation 22:20). Christian hope rests on the fidelity of the person to whom we pray. We can still ask him to come after we have received him in Holy Communion. His sacramental presence is real. We hunger for the day when the veil is lifted and those who have slept in him rise.

Christians learned to wait from Israel. When Isaiah sees the nations going up to Jerusalem, they go there to be taught. The instruction they receive makes a difference to what a smith does with his iron (Isaiah 2:1--5). When the Virgin of Nazareth sings about the mercy shown to Abraham, she remembers her own people and a promise older than her household (Luke 1:54--55). At the beginning of the liturgical year we sit down among the people who first heard those promises. Each December the Church gives her children time to learn their words. The \emph{Catechism of the Catholic Church}, no.~524, explains how sharing Israel's long expectation rekindles our desire for the Lord's return.

In the Creed we confess that the same Jesus who was born of Mary will come to judge the living and the dead. A judge who has himself been hungry, dependent on a mother's care, betrayed by a friend and buried by loving hands knows the people who will stand before him. Let that truth reach the conscience. If someone is waiting for wages you owe or for you to tell the truth about an injury, go to that person. The First Sunday's prayer asks God to give us the will for just such deeds. Grace reaches a person's hands.
